
\documentclass{ximera}
\title{Inverses CW3}

\newcommand{\pskip}{\vskip 0.1 in}

\begin{document}
\begin{abstract}
Review of inverse functions.
\end{abstract}
\maketitle

\begin{question}  \label{Qdf3r344xxz}
The function
\[
       Q = f(P) \, , \, 5 \leq P \leq 15,
\]
expresses the average number of burgers a diner sells per day in terms of the price (in dollars/burger).

\begin{enumerate}

\item Would you expect $f$ to be a one-to-one function? Explain your reasoning.

\item Interpret the meanings of each of the following expressions.

\begin{enumerate}
\item $f(10)$

\item $f^{-1}(30)$.

\item $f^{-1}(f(6)-10)$

\item $f(f^{-1}(30)-2)$

%\item $f^{-1}(f(P)-10)$

%\item $Q - f(f^{-1}(Q)-2)$

\end{enumerate}

\item Suppose the function $f$ is linear. Suppose also the diner sells an average of $50$ burgers/day when the price is $\$5$/burger and an average of $10$ burgers/day when the price is $\$15$/burger.

\begin{enumerate}
\item Find expressions for the functions $Q=f(P)$ and $P=f^{-1}(Q)$. Include a domain and range for 	each.

\item Evaluate $f(f^{-1}(30)-2)$

\item Let 
\[
      q = g(Q) = f(f^{-1}(Q)-2).   
\]

\begin{enumerate}
\item Explain the meaning of the function $g$. What does it take as an input? What does it return as an output?

\item Use algebra to find a simplified expression for the function $q=g(Q)$.

\item Try to use logic instead of algebra to find a simplified expression for the function $q=g(Q)$.

\item Use set builder notation to write the domain and range of the function $q=g(Q)$.

\item Graph the function $q=g(Q)$. Label the coordinate axes with the appropriate variable names and units.
\end{enumerate}
\end{enumerate}
\end{enumerate}

\end{question}

\begin{question} \label{Qpve03}
The function
\[
      V = f(t) \, , \, 0\leq t \leq 6,
\]
expresses the volume of water in a tank (measured in gallons) in terms of the number of minutes past noon.

\pdfOnly{
Access Desmos interactives through the online version of this text at
 

\begin{enumerate}
\item Explain what it would mean for the function $f$ to be one-to-one in this particular scenario. Use everyday English. Avoid math terms like input and output.

\item Explain what it would mean for the function $f$ to \emph{not} be one-to-one in this particular scenario. Use everyday English. Avoid math terms like input and output.  

\item Now suppose 
\[
  V = f(t) = \frac{4}{3}\left( t-6 \right)^2 \, , \, 0\leq t \leq 6.
\]


\begin{enumerate}
\item Sketch by hand a graph of the function $f$. Label the axes with the appropriate variable names and units. 

\item Use your graph to determine if the function $f$ is one-to-one. Explain your reasoning in the context of this particular scenario.

\item Describe the action of the function $f^{-1}$ in this particular scenario. What does it take as an input? What does it return as an output.

\item Find an expression for the function $t = f^{-1}(V)$. 

\item Write the domain and range of $f^{-1}$ in set notation.

\item Simplify the composition $f\circ f^{-1}$ algebraically. Use the appropriate variable name for the input.

\item Simplify the composition $f^{-1}\circ f$ algebraically. Use the appropriate variable name for the input.

\end{enumerate}
\end{enumerate}
\end{question}

\begin{question}\label{Qpfef3z}
The function
\[
  h = f(t) = 320 - 5(t-6)^2 \, , \, 0\leq t \leq 14.
\]
expresses the height (in meters) of a rock in terms of the number of seconds since it was launched straight up. Its graph is shown below.

\href{https://www.desmos.com/calculator/ew9jsipj5z}.
}
 
\begin{onlineOnly}
    \begin{center}
\desmos{nho1omvdtd}{900}{600}
\end{center}
\end{onlineOnly}

\href{https://www.desmos.com/calculator/ew9jsipj5z}{141:Rock}

\begin{enumerate}
\item Is function $f$ one-to-one? Explain your reasoning in this particular scenario.

\item Find a function $T=u(h)$ that takes as an input a height (in meters) and returns as an output the time (in seconds past noon) when the rock was at that height and on its way up. Include the domain.

\item Explain the meaning of the composition $u\circ f$.

\item Find the domain of the composition $u\circ f$. Write it as a set with the appropriate variable.

\item Use the graph of the function $f$ above to sketch a graph of the function $T= u(f(t))$. Explain your reasoning.

\item Simplify the composition $T=u(f(t))$ alegebraically. Your final expression should involve the absolute value function. Then graph the composition and see how it compares with your sketch from part (d).  %Then find a piecewise-defined equivalent form of the composition.

\item Find the inverse of the function $u$. Stop and think. You need not do any algebra for this.

\item Repeat parts (a)-(f) for the function $T=d(h)$ that takes as an input a height (in meters) and returns as an output the time (in seconds past noon) when the rock was at that height and on its way down.
\end{enumerate}
\end{question}

\end{document}